\documentclass[11pt]{article}
\usepackage{amsmath, amssymb, geometry, graphicx}
\geometry{a4paper, margin=1in}

\title{Practice Problem Set: Double Integrals in Polar Coordinates}
\author{Homework 15.3 Extension}
\date{}

\begin{document}

\maketitle

\section*{Concept Checklist}
Based on the foundational material and advanced diagnostic questions, the following concepts are critical for mastering double integrals in polar coordinates. These concepts are rigorously tested in the problem set below.

\begin{itemize}
    \item \textbf{C1: Cartesian to Polar Region Conversion \& Setup:} Identifying the correct radial bounds ($dr$) and angular bounds ($d\theta$) from graphs, Cartesian boundaries, or text descriptions.
    \item \textbf{C2: Evaluating Polar Integrals over Disks and Sectors:} Calculating integrals over full circles, half-circles, and quadrant slices.
    \item \textbf{C3: Evaluating Polar Integrals over Annuli (Rings):} Setting up bounds with an inner and outer radius ($a \le r \le b$).
    \item \textbf{C4: Regions Bounded by Shifted Polar Curves:} Setting up bounds for shifted circles like $r = a \cos\theta$ or $r = a \sin\theta$.
    \item \textbf{C5: U-Substitution with the Jacobian:} Recognizing that the conversion $dA = r \, dr \, d\theta$ perfectly sets up integrals involving $e^{-r^2}$ or trigonometric functions of $r^2$.
    \item \textbf{C6: Trigonometric Integration \& Half-Angle Power Reduction:} Using identities like $\sin^2\theta = \frac{1-\cos(2\theta)}{2}$ to solve angular integrals.
    \item \textbf{C7: Volume Under One Surface:} Finding the 3D volume trapped between a function $z = f(x,y)$ and the $xy$-plane.
    \item \textbf{C8: Volume Between Two Surfaces:} Equating two functions to find their intersection in the $xy$-plane, and integrating $(Top - Bottom)$.
    \item \textbf{C9: Real-World Applications \& Modeling:} Applying polar integrals to center of mass, variable density, and Gaussian probability distributions.
    \item \textbf{C10: Diagnostic Testing (Finding the Flaw):} Identifying common mistakes such as forgetting the Jacobian '$r$', applying wrong angular bounds, or squaring algebraic terms incorrectly.
\end{itemize}

\newpage
\section*{30 Practice Problems and Solutions}

\textbf{Problem 1:} Set up, but do not evaluate, the polar integral for $\iint_D f(x,y) \, dA$, where $D$ is the sector of a circle of radius 6 located strictly in the third quadrant.
\\
\textbf{Solution:} 
The third quadrant restricts the angle to $\pi \le \theta \le \frac{3\pi}{2}$. The radius spans from the origin to 6, so $0 \le r \le 6$. 
Substitute $x = r\cos\theta$, $y = r\sin\theta$, and $dA = r \, dr \, d\theta$.
\\
\textbf{Final Answer:} $\int_{\pi}^{3\pi/2} \int_0^6 f(r\cos\theta, r\sin\theta) \, r \, dr \, d\theta$

\vspace{0.5cm}
\textbf{Problem 2:} Set up, but do not evaluate, the polar integral for $\iint_D f(x,y) \, dA$, where $D$ is the right half of the annulus bounded by $x^2+y^2=9$ and $x^2+y^2=25$.
\\
\textbf{Solution:}
The region is between radii 3 and 5, so $3 \le r \le 5$. The "right half" means $x \ge 0$, which corresponds to the first and fourth quadrants. The angular bounds are $-\pi/2 \le \theta \le \pi/2$.
\\
\textbf{Final Answer:} $\int_{-\pi/2}^{\pi/2} \int_3^5 f(r\cos\theta, r\sin\theta) \, r \, dr \, d\theta$

\vspace{0.5cm}
\textbf{Problem 3:} Evaluate $\iint_D 2xy^2 \, dA$, where $D$ is the upper half of the disk $x^2+y^2 \le 9$.
\\
\textbf{Solution:}
Upper half disk means $0 \le \theta \le \pi$ and $0 \le r \le 3$. 
Integrand: $2(r\cos\theta)(r\sin\theta)^2 \cdot r = 2r^4\cos\theta\sin^2\theta$.
$$ \int_0^\pi \int_0^3 2r^4\cos\theta\sin^2\theta \, dr \, d\theta = \left( \int_0^\pi \cos\theta\sin^2\theta \, d\theta \right) \left( \int_0^3 2r^4 \, dr \right) $$
For the $\theta$-integral, let $u = \sin\theta \implies du = \cos\theta \, d\theta$. The bounds go from $\sin(0)=0$ to $\sin(\pi)=0$. Since the bounds match, the integral evaluates to 0. 
\\
\textbf{Final Answer:} $0$

\vspace{0.5cm}
\textbf{Problem 4:} Find the volume of the solid bounded above by the paraboloid $z = x^2+y^2$ and below by the region in the $xy$-plane bounded by $y=x$, the $x$-axis, and the circle $x^2+y^2=4$.
\\
\textbf{Solution:}
The region in the $xy$-plane is bounded by $\theta = 0$ (x-axis) and $\theta = \pi/4$ (the line $y=x$), with radius $0 \le r \le 2$. The top surface is $z = r^2$.
$$ V = \int_0^{\pi/4} \int_0^2 (r^2) \, r \, dr \, d\theta = \int_0^{\pi/4} d\theta \int_0^2 r^3 \, dr = \left( \frac{\pi}{4} \right) \left[ \frac{r^4}{4} \right]_0^2 = \frac{\pi}{4} \left( \frac{16}{4} \right) = \pi $$
\textbf{Final Answer:} $\pi$

\vspace{0.5cm}
\textbf{Problem 5:} Evaluate the iterated integral by first converting to polar coordinates: $\int_0^3 \int_{-\sqrt{9-y^2}}^{\sqrt{9-y^2}} (x^2+y^2)^{3/2} \, dx \, dy$.
\\
\textbf{Solution:}
The limits for $x$ are $-\sqrt{9-y^2}$ to $\sqrt{9-y^2}$, forming a full horizontal span of a circle. However, $y$ only goes from 0 to 3. This outlines the upper half of a disk of radius 3.
Bounds: $0 \le \theta \le \pi$, $0 \le r \le 3$. 
Integrand: $(r^2)^{3/2} \cdot r = r^3 \cdot r = r^4$.
$$ \int_0^\pi \int_0^3 r^4 \, dr \, d\theta = \pi \left[ \frac{r^5}{5} \right]_0^3 = \frac{243\pi}{5} $$
\textbf{Final Answer:} $\frac{243\pi}{5}$

\vspace{0.5cm}
\textbf{Problem 6:} Evaluate $\iint_D e^{-(x^2+y^2)} \, dA$, where $D$ is the disk of radius 2 centered at the origin.
\\
\textbf{Solution:}
Bounds: $0 \le \theta \le 2\pi$, $0 \le r \le 2$.
Integrand: $e^{-r^2} \, r$.
$$ \int_0^{2\pi} \int_0^2 e^{-r^2} r \, dr \, d\theta = 2\pi \int_0^2 e^{-r^2} r \, dr $$
Let $u = -r^2$, $du = -2r \, dr \implies -\frac{1}{2}du = r \, dr$.
$$ 2\pi \left[ -\frac{1}{2} e^{-r^2} \right]_0^2 = -\pi (e^{-4} - 1) = \pi(1 - e^{-4}) $$
\textbf{Final Answer:} $\pi(1 - e^{-4})$

\vspace{0.5cm}
\textbf{Problem 7:} Evaluate $\iint_D \sin(x^2+y^2) \, dA$, where $D$ is the annulus $\frac{\sqrt{\pi}}{2} \le \sqrt{x^2+y^2} \le \sqrt{\pi}$.
\\
\textbf{Solution:}
Bounds: $0 \le \theta \le 2\pi$, $\frac{\sqrt{\pi}}{2} \le r \le \sqrt{\pi}$.
Integrand: $\sin(r^2) \, r$.
$$ \int_0^{2\pi} \int_{\sqrt{\pi}/2}^{\sqrt{\pi}} \sin(r^2) r \, dr \, d\theta = 2\pi \left[ -\frac{1}{2} \cos(r^2) \right]_{\sqrt{\pi}/2}^{\sqrt{\pi}} $$
$$ = -\pi \left( \cos(\pi) - \cos\left(\frac{\pi}{4}\right) \right) = -\pi \left( -1 - \frac{\sqrt{2}}{2} \right) = \pi \left( 1 + \frac{\sqrt{2}}{2} \right) $$
\textbf{Final Answer:} $\pi \left( 1 + \frac{\sqrt{2}}{2} \right)$

\vspace{0.5cm}
\textbf{Problem 8:} Find the area of the region bounded by the circle $r = 4\sin\theta$.
\\
\textbf{Solution:}
The curve $r = 4\sin\theta$ traces a full circle in the upper half-plane as $\theta$ ranges from $0$ to $\pi$. Area is $\iint 1 \, dA$.
$$ A = \int_0^\pi \int_0^{4\sin\theta} r \, dr \, d\theta = \int_0^\pi \left[ \frac{r^2}{2} \right]_0^{4\sin\theta} d\theta = \int_0^\pi 8\sin^2\theta \, d\theta $$
Use the half-angle identity: $8 \int_0^\pi \frac{1-\cos(2\theta)}{2} d\theta = 4 \left[ \theta - \frac{\sin(2\theta)}{2} \right]_0^\pi = 4(\pi - 0) = 4\pi $.
\\
\textbf{Final Answer:} $4\pi$

\vspace{0.5cm}
\textbf{Problem 9:} Evaluate $\iint_D \cos\left(\sqrt{x^2+y^2}\right) \, dA$ where $D$ is the disk with center at the origin and radius 2.
\\
\textbf{Solution:}
Bounds: $0 \le \theta \le 2\pi$, $0 \le r \le 2$. Integrand: $\cos(r) \, r$.
$$ \int_0^{2\pi} \int_0^2 r\cos(r) \, dr \, d\theta = 2\pi \int_0^2 r\cos(r) \, dr $$
Use integration by parts: let $u=r$, $dv=\cos(r)dr \implies du=dr$, $v=\sin(r)$.
$$ \int r\cos(r) \, dr = r\sin(r) - \int \sin(r) \, dr = r\sin(r) + \cos(r) $$
Evaluate from 0 to 2: $(2\sin(2) + \cos(2)) - (0 + \cos(0)) = 2\sin(2) + \cos(2) - 1$. Multiply by $2\pi$.
\\
\textbf{Final Answer:} $2\pi(2\sin(2) + \cos(2) - 1)$

\vspace{0.5cm}
\textbf{Problem 10:} Find the volume of the solid bounded above by the paraboloid $z = 16 - x^2 - y^2$ and below by the $xy$-plane ($z=0$).
\\
\textbf{Solution:}
Find intersection with $z=0$ to get bounds: $16 - x^2 - y^2 = 0 \implies r^2 = 16 \implies r = 4$.
Bounds: $0 \le \theta \le 2\pi$, $0 \le r \le 4$.
$$ V = \int_0^{2\pi} \int_0^4 (16 - r^2) r \, dr \, d\theta = 2\pi \int_0^4 (16r - r^3) \, dr $$
$$ = 2\pi \left[ 8r^2 - \frac{r^4}{4} \right]_0^4 = 2\pi \left( 8(16) - \frac{256}{4} \right) = 2\pi (128 - 64) = 128\pi $$
\textbf{Final Answer:} $128\pi$

\vspace{0.5cm}
\textbf{Problem 11:} Find the volume under the surface $z = xy$ and above the first quadrant of the disk $x^2+y^2 \le 16$.
\\
\textbf{Solution:}
Bounds: $0 \le \theta \le \pi/2$, $0 \le r \le 4$.
Integrand: $(r\cos\theta)(r\sin\theta) \cdot r = r^3 \cos\theta \sin\theta$.
$$ \int_0^{\pi/2} \int_0^4 r^3 \cos\theta\sin\theta \, dr \, d\theta = \left( \int_0^{\pi/2} \cos\theta\sin\theta \, d\theta \right) \left( \int_0^4 r^3 \, dr \right) $$
$\theta$-integral: $\left[ \frac{\sin^2\theta}{2} \right]_0^{\pi/2} = \frac{1}{2} - 0 = \frac{1}{2}$.
$r$-integral: $\left[ \frac{r^4}{4} \right]_0^4 = 64$.
Multiply: $\frac{1}{2} \times 64 = 32$.
\\
\textbf{Final Answer:} $32$

\vspace{0.5cm}
\textbf{Problem 12:} Set up and evaluate the integral to find the area of the annular sector between radii 1 and 3, bound between $\theta = \pi/4$ and $\theta = 3\pi/4$.
\\
\textbf{Solution:}
Area $A = \iint_D 1 \, dA = \int_{\pi/4}^{3\pi/4} \int_1^3 r \, dr \, d\theta $.
$$ = \left( \frac{3\pi}{4} - \frac{\pi}{4} \right) \left[ \frac{r^2}{2} \right]_1^3 = \left( \frac{2\pi}{4} \right) \left( \frac{9}{2} - \frac{1}{2} \right) = \frac{\pi}{2} (4) = 2\pi $$
\textbf{Final Answer:} $2\pi$

\vspace{0.5cm}
\textbf{Problem 13:} Find the volume of the solid bounded by the paraboloid $z = 3x^2+3y^2$ and the plane $z = 27$.
\\
\textbf{Solution:}
Top surface: $z=27$. Bottom surface: $z=3r^2$.
Intersection: $27 = 3r^2 \implies r^2 = 9 \implies r = 3$.
Bounds: $0 \le \theta \le 2\pi$, $0 \le r \le 3$.
$$ V = \int_0^{2\pi} \int_0^3 (27 - 3r^2)r \, dr \, d\theta = 2\pi \int_0^3 (27r - 3r^3) \, dr $$
$$ = 2\pi \left[ \frac{27}{2}r^2 - \frac{3}{4}r^4 \right]_0^3 = 2\pi \left( \frac{243}{2} - \frac{243}{4} \right) = 2\pi \left( \frac{243}{4} \right) = \frac{243\pi}{2} $$
\textbf{Final Answer:} $\frac{243\pi}{2}$

\vspace{0.5cm}
\textbf{Problem 14:} Find the volume of the solid bounded above by $z = 20 - x^2 - y^2$ and below by $z = x^2+y^2$.
\\
\textbf{Solution:}
Intersection: $20 - r^2 = r^2 \implies 2r^2 = 20 \implies r^2 = 10 \implies r = \sqrt{10}$.
$$ V = \int_0^{2\pi} \int_0^{\sqrt{10}} \left( (20 - r^2) - r^2 \right) r \, dr \, d\theta = 2\pi \int_0^{\sqrt{10}} (20r - 2r^3) \, dr $$
$$ = 2\pi \left[ 10r^2 - \frac{2r^4}{4} \right]_0^{\sqrt{10}} = 2\pi \left( 10(10) - \frac{1}{2}(100) \right) = 2\pi (100 - 50) = 100\pi $$
\textbf{Final Answer:} $100\pi$

\vspace{0.5cm}
\textbf{Problem 15:} Find the volume of the solid inside the sphere $x^2+y^2+z^2=16$ and outside the cylinder $x^2+y^2=4$.
\\
\textbf{Solution:}
Utilizing symmetry, we can calculate the volume of the top half and multiply by 2.
Top surface: $z = \sqrt{16-r^2}$. The region $D$ is an annulus from the cylinder wall $r=2$ to the sphere's maximum radius in the $xy$-plane $r=4$.
$$ V_{top} = \int_0^{2\pi} \int_2^4 \sqrt{16-r^2} \, r \, dr \, d\theta = 2\pi \int_2^4 \sqrt{16-r^2} \, r \, dr $$
Let $u = 16-r^2$, $du = -2r \, dr \implies -\frac{1}{2}du = r \, dr$. 
$$ = 2\pi \left[ -\frac{1}{3}(16-r^2)^{3/2} \right]_2^4 = 2\pi \left( 0 - \left(-\frac{1}{3}(12)^{3/2}\right) \right) = \frac{2\pi}{3} (12\sqrt{12}) = \frac{2\pi}{3} (24\sqrt{3}) = 16\pi\sqrt{3} $$
Total Volume = $2 \times V_{top} = 32\pi\sqrt{3}$.
\\
\textbf{Final Answer:} $32\pi\sqrt{3}$

\vspace{0.5cm}
\textbf{Problem 16:} Find the volume of the solid above the cone $z = \sqrt{x^2+y^2}$ and below the sphere $x^2+y^2+z^2=8$.
\\
\textbf{Solution:}
Cone: $z=r$. Sphere: $z = \sqrt{8-r^2}$.
Intersection: $r = \sqrt{8-r^2} \implies r^2 = 8-r^2 \implies 2r^2=8 \implies r=2$.
$$ V = \int_0^{2\pi} \int_0^2 \left( \sqrt{8-r^2} - r \right) r \, dr \, d\theta = 2\pi \left( \int_0^2 r\sqrt{8-r^2} \, dr - \int_0^2 r^2 \, dr \right) $$
For the first integral, let $u=8-r^2$.
$$ 2\pi \left( \left[ -\frac{1}{3}(8-r^2)^{3/2} \right]_0^2 - \left[ \frac{r^3}{3} \right]_0^2 \right) = 2\pi \left( \left(-\frac{1}{3}(4)^{3/2} + \frac{1}{3}(8)^{3/2} \right) - \frac{8}{3} \right) $$
$$ = 2\pi \left( -\frac{8}{3} + \frac{16\sqrt{2}}{3} - \frac{8}{3} \right) = 2\pi \left( \frac{16\sqrt{2} - 16}{3} \right) = \frac{32\pi(\sqrt{2}-1)}{3} $$
\textbf{Final Answer:} $\frac{32\pi(\sqrt{2}-1)}{3}$

\vspace{0.5cm}
\textbf{Problem 17:} Evaluate $\iint_D 4xy \, dA$, where $D$ is the region in the first quadrant bounded by the circles $x^2+y^2=4$ and $x^2+y^2=25$.
\\
\textbf{Solution:}
Bounds: $0 \le \theta \le \pi/2$, $2 \le r \le 5$. Integrand: $4(r\cos\theta)(r\sin\theta) \cdot r = 4r^3\cos\theta\sin\theta$.
$$ \int_0^{\pi/2} 4\cos\theta\sin\theta \, d\theta \times \int_2^5 r^3 \, dr = \int_0^{\pi/2} 2\sin(2\theta) \, d\theta \times \left[ \frac{r^4}{4} \right]_2^5 $$
$$ = \left[ -\cos(2\theta) \right]_0^{\pi/2} \times \left( \frac{625}{4} - \frac{16}{4} \right) = (1 - (-1)) \times \frac{609}{4} = 2 \times \frac{609}{4} = \frac{609}{2} $$
\textbf{Final Answer:} $\frac{609}{2}$

\vspace{0.5cm}
\textbf{Problem 18:} Convert to polar coordinates and evaluate: $\int_0^{\sqrt{2}} \int_y^{\sqrt{4-y^2}} \frac{1}{1+x^2+y^2} \, dx \, dy$.
\\
\textbf{Solution:}
The region has $y$ ranging from $0$ to $\sqrt{2}$. For a given $y$, $x$ ranges from $y$ (the line $y=x$) to $\sqrt{4-y^2}$ (the right half of the circle $r=2$).
In polar coordinates, the line $y=x$ corresponds to $\theta = \pi/4$, and the $x$-axis is $\theta = 0$. Since $y \ge 0$ and $x \ge y$, the region is bounded by $0 \le \theta \le \pi/4$ and $0 \le r \le 2$.
$$ \int_0^{\pi/4} \int_0^2 \frac{1}{1+r^2} \, r \, dr \, d\theta = \frac{\pi}{4} \left[ \frac{1}{2} \ln(1+r^2) \right]_0^2 = \frac{\pi}{8} \ln(5) $$
\textbf{Final Answer:} $\frac{\pi}{8} \ln(5)$

\vspace{0.5cm}
\textbf{Problem 19: Diagnostic Testing - Find the Flaw.} A student evaluates $\iint_D x^2 \, dA$, where $D$ is the left half of the disk $r \le 3$. Their setup is: $\int_{\pi/2}^{3\pi/2} \int_0^3 r^2 \cos^2\theta \, dr \, d\theta$. Explain the error and correct it.
\\
\textbf{Solution:}
\textit{Flaw:} The student forgot to multiply by the Jacobian $r$ from $dA = r \, dr \, d\theta$. While $x^2 = r^2\cos^2\theta$, the area differential adds an extra $r$.
\textit{Correction:} The integrand should be $r^3 \cos^2\theta$. The correct integral is $\int_{\pi/2}^{3\pi/2} \int_0^3 r^3 \cos^2\theta \, dr \, d\theta$.

\vspace{0.5cm}
\textbf{Problem 20:} Find the mass of a circular plate of radius 2 centered at the origin if its density function is $\rho(x,y) = e^{-(x^2+y^2)}$.
\\
\textbf{Solution:}
Mass $m = \iint_D \rho(x,y) \, dA$. Bounds: $0 \le \theta \le 2\pi$, $0 \le r \le 2$.
$$ m = \int_0^{2\pi} \int_0^2 e^{-r^2} \, r \, dr \, d\theta = 2\pi \left[ -\frac{1}{2} e^{-r^2} \right]_0^2 = \pi(1 - e^{-4}) $$
\textbf{Final Answer:} $\pi(1 - e^{-4})$

\vspace{0.5cm}
\textbf{Problem 21:} Convert to polar and evaluate: $\int_{-3}^3 \int_0^{\sqrt{9-x^2}} \sin(x^2+y^2) \, dy \, dx$.
\\
\textbf{Solution:}
$x$ ranges from $-3$ to $3$, and $y$ ranges from $0$ to $\sqrt{9-x^2}$. This represents the upper half of a disk of radius 3. Bounds: $0 \le \theta \le \pi$, $0 \le r \le 3$.
$$ \int_0^\pi \int_0^3 \sin(r^2) \, r \, dr \, d\theta = \pi \left[ -\frac{1}{2} \cos(r^2) \right]_0^3 = -\frac{\pi}{2} (\cos(9) - \cos(0)) = \frac{\pi}{2}(1 - \cos(9)) $$
\textbf{Final Answer:} $\frac{\pi}{2}(1 - \cos(9))$

\vspace{0.5cm}
\textbf{Problem 22:} Evaluate $\iint_D \frac{y^2}{x^2+y^2} \, dA$, where $D$ is the region between $r=2$ and $r=4$.
\\
\textbf{Solution:}
Bounds: $0 \le \theta \le 2\pi$, $2 \le r \le 4$.
Integrand: $\frac{r^2\sin^2\theta}{r^2} \cdot r = r\sin^2\theta$.
$$ \int_0^{2\pi} \int_2^4 r\sin^2\theta \, dr \, d\theta = \int_0^{2\pi} \sin^2\theta \, d\theta \times \int_2^4 r \, dr = \left( \pi \right) \left[ \frac{r^2}{2} \right]_2^4 = \pi(8 - 2) = 6\pi $$
\textbf{Final Answer:} $6\pi$

\vspace{0.5cm}
\textbf{Problem 23:} Find the volume of the solid bounded by the paraboloid $z = 9 - x^2 - y^2$ and the $xy$-plane, strictly in the first octant ($x \ge 0, y \ge 0$).
\\
\textbf{Solution:}
Bounds: $0 \le \theta \le \pi/2$ (first octant restricts to Q1 in $xy$-plane). Base radius: $9 - r^2 = 0 \implies r = 3$.
$$ V = \int_0^{\pi/2} \int_0^3 (9 - r^2) r \, dr \, d\theta = \frac{\pi}{2} \left[ \frac{9}{2}r^2 - \frac{r^4}{4} \right]_0^3 = \frac{\pi}{2} \left( \frac{81}{2} - \frac{81}{4} \right) = \frac{\pi}{2} \left( \frac{81}{4} \right) = \frac{81\pi}{8} $$
\textbf{Final Answer:} $\frac{81\pi}{8}$

\vspace{0.5cm}
\textbf{Problem 24:} Find the volume under $z = x^2+y^2$ and above the circular region bounded by $r = 2\sin\theta$.
\\
\textbf{Solution:}
Bounds: $0 \le \theta \le \pi$, $0 \le r \le 2\sin\theta$.
$$ V = \int_0^\pi \int_0^{2\sin\theta} (r^2) \, r \, dr \, d\theta = \int_0^\pi \left[ \frac{r^4}{4} \right]_0^{2\sin\theta} d\theta = \int_0^\pi 4\sin^4\theta \, d\theta $$
Use power reduction: $\sin^4\theta = \left(\frac{1-\cos(2\theta)}{2}\right)^2 = \frac{1 - 2\cos(2\theta) + \cos^2(2\theta)}{4} = \frac{1 - 2\cos(2\theta) + \frac{1+\cos(4\theta)}{2}}{4} = \frac{3/2 - 2\cos(2\theta) + 1/2\cos(4\theta)}{4}$. 
Multiplying by 4: $\int_0^\pi \left( \frac{3}{2} - 2\cos(2\theta) + \frac{1}{2}\cos(4\theta) \right) d\theta$.
The cosine terms integrate to 0 over $[0, \pi]$.
$V = \left[ \frac{3}{2}\theta \right]_0^\pi = \frac{3\pi}{2} $.
\\
\textbf{Final Answer:} $\frac{3\pi}{2}$

\vspace{0.5cm}
\textbf{Problem 25:} A financial quant models the joint probability of two assets using $p(x,y) = \frac{1}{8\pi} e^{-(x^2+y^2)/8}$. Find the probability that the combined return falls within a radius of $R=4$ from the mean.
\\
\textbf{Solution:}
$P = \iint_D \frac{1}{8\pi} e^{-(x^2+y^2)/8} \, dA$ over $r \le 4$.
$$ P = \int_0^{2\pi} \int_0^4 \frac{1}{8\pi} e^{-r^2/8} \, r \, dr \, d\theta = 2\pi \left( \frac{1}{8\pi} \right) \int_0^4 e^{-r^2/8} r \, dr = \frac{1}{4} \int_0^4 e^{-r^2/8} r \, dr $$
Let $u = -r^2/8$, $du = -r/4 \, dr$. When $r=0$, $u=0$. When $r=4$, $u=-2$.
$$ \int_0^{-2} -e^u \, du = \left[ -e^u \right]_0^{-2} = -e^{-2} - (-1) = 1 - e^{-2} $$
\textbf{Final Answer:} $1 - e^{-2}$

\vspace{0.5cm}
\textbf{Problem 26: Diagnostic Testing - Find the Flaw.} To find the volume bounded by $z=10-x^2-y^2$ and $z=2x^2+2y^2$, a student calculates intersection $r=\sqrt{10/3}$, and evaluates $\int_0^{2\pi} \int_0^{\sqrt{10/3}} \left(2r^2 - (10-r^2)\right) r \, dr \, d\theta$, resulting in a negative volume.
\\
\textbf{Solution:}
\textit{Flaw:} The student subtracted the top function from the bottom function (Floor - Ceiling). At $r=0$, $10-0^2=10$ (top), while $2(0)^2=0$ (bottom).
\textit{Correction:} The integrand should be $(10-r^2) - 2r^2 = 10-3r^2$. The correct setup is $\int_0^{2\pi} \int_0^{\sqrt{10/3}} (10-3r^2)r \, dr \, d\theta$.

\vspace{0.5cm}
\textbf{Problem 27:} Find the mass of a lamina occupying the first quadrant of the annulus $1 \le x^2+y^2 \le 9$ if the density is $\rho(x,y) = \frac{1}{x^2+y^2}$.
\\
\textbf{Solution:}
Bounds: $0 \le \theta \le \pi/2$, $1 \le r \le 3$.
Integrand: $\frac{1}{r^2} \cdot r = \frac{1}{r}$.
$$ m = \int_0^{\pi/2} \int_1^3 \frac{1}{r} \, dr \, d\theta = \frac{\pi}{2} \left[ \ln(r) \right]_1^3 = \frac{\pi}{2}(\ln(3) - \ln(1)) = \frac{\pi}{2}\ln(3) $$
\textbf{Final Answer:} $\frac{\pi\ln(3)}{2}$

\vspace{0.5cm}
\textbf{Problem 28:} Find the volume of the solid bounded above by the hemisphere $z = \sqrt{16-x^2-y^2}$ and below by the $xy$-plane.
\\
\textbf{Solution:}
The base is the full circle of radius 4. Bounds: $0 \le \theta \le 2\pi$, $0 \le r \le 4$.
$$ V = \int_0^{2\pi} \int_0^4 \sqrt{16-r^2} \, r \, dr \, d\theta = 2\pi \left[ -\frac{1}{3}(16-r^2)^{3/2} \right]_0^4 = 2\pi \left( 0 - \left(-\frac{1}{3}(16)^{3/2}\right) \right) = 2\pi \left( \frac{64}{3} \right) = \frac{128\pi}{3} $$
\textbf{Final Answer:} $\frac{128\pi}{3}$

\vspace{0.5cm}
\textbf{Problem 29:} Convert and evaluate: $\int_0^5 \int_0^{\sqrt{25-y^2}} 3 \, dx \, dy$.
\\
\textbf{Solution:}
$y$ ranges from 0 to 5, and $x$ ranges from 0 to $\sqrt{25-y^2}$. This represents the first quadrant of a disk of radius 5.
Bounds: $0 \le \theta \le \pi/2$, $0 \le r \le 5$.
$$ \int_0^{\pi/2} \int_0^5 3r \, dr \, d\theta = \frac{\pi}{2} \left[ \frac{3r^2}{2} \right]_0^5 = \frac{\pi}{2} \left( \frac{75}{2} \right) = \frac{75\pi}{4} $$
\textbf{Final Answer:} $\frac{75\pi}{4}$

\vspace{0.5cm}
\textbf{Problem 30: Diagnostic Testing - Find the Flaw.} A student converts the integral $\iint_D (x^2+y^2)^2 \, dA$ into polar as $\iint_D r^2 \cdot r \, dr \, d\theta = \iint_D r^3 \, dr \, d\theta$.
\\
\textbf{Solution:}
\textit{Flaw:} The student failed to square the polar substitution properly. $(x^2+y^2)^2$ translates to $(r^2)^2 = r^4$. 
\textit{Correction:} Multiplying by the Jacobian $r$, the correct integrand is $r^4 \cdot r = r^5$. The proper integral is $\iint_D r^5 \, dr \, d\theta$.

\newpage
\section*{Mapping of Concepts to Problem Numbers}
Use this key to cross-reference areas where you need additional practice with the specific problems that target them:

\begin{itemize}
    \item \textbf{C1: Cartesian to Polar Region Conversion \& Setup:} 1, 2, 5, 18, 21, 29
    \item \textbf{C2: Evaluating Polar Integrals over Disks and Sectors:} 3, 4, 5, 11, 18, 29
    \item \textbf{C3: Evaluating Polar Integrals over Annuli (Rings):} 7, 12, 17, 22, 27
    \item \textbf{C4: Regions Bounded by Shifted Polar Curves:} 8, 24
    \item \textbf{C5: U-Substitution with the Jacobian:} 6, 7, 9, 20, 21, 25
    \item \textbf{C6: Trigonometric Integration \& Half-Angle Power Reduction:} 3, 8, 11, 17, 22, 24
    \item \textbf{C7: Volume Under One Surface:} 4, 10, 23, 24, 28
    \item \textbf{C8: Volume Between Two Surfaces \& Finding Intersections:} 13, 14, 15, 16, 26
    \item \textbf{C9: Applications: Variable Density, Center of Mass, Probability:} 20, 25, 27
    \item \textbf{C10: Diagnostic Testing: Identifying Common Setup Flaws:} 19, 26, 30
\end{itemize}

\end{document}